% ============================================================
\section{Conclusion}
\label{sec:conclusion}
% ============================================================

In this paper, we present \textsc{DeepTutor}, an open-source framework for personalized tutoring built around a shared personalization substrate and organized along two complementary dimensions.

The first dimension, \emph{personalized agentic tutoring}, establishes the validated core of the system.
A hybrid personalization engine couples static knowledge grounding with a dynamic trace forest to close a bidirectional loop between citation-grounded problem solving and difficulty-calibrated question generation.
The same engine extends to collaborative writing, deep research, and guided learning, ensuring that progress in any modality informs the others.
The second dimension, \emph{proactive autonomous companionship}, deploys these capabilities in a broader setting through TutorBot, where autonomous agents operate through extensible skills, multi-bot coordination, and context-unified multi-channel access.

Evaluation of the foundational pipeline on TutorBench, together with a first-person interactive protocol designed from the learner's perspective, demonstrates a 10.8\% improvement on personalized tutoring metrics and an average 28.6\% gain on general reasoning benchmarks across five backbone models.
Ablation studies confirm that knowledge grounding and learner memory contribute along largely complementary axes, validating the hybrid design upon which the broader system is built.

\paragraph{Broader Implications.}
The architectural pattern explored here extends beyond education.
Whenever an AI system must personalize over time while remaining coherent across tools and interfaces, a shared personalization substrate paired with a proactive deployment layer offers a useful design template.
\textsc{DeepTutor} can be viewed as one concrete instantiation of this broader class of agentic personalized systems.

\paragraph{Future Work.}
Longitudinal deployment studies with real students would provide the most rigorous validation of the cognitive extensions and the proactive multi-agent layer.
Collaborative writing, deep research, guided learning, and cross-channel TutorBot behaviors should be evaluated with deployment-scale metrics that capture retention, follow-up quality, and sustained learner engagement.
The trace forest could benefit from more sophisticated consolidation strategies, such as importance-weighted compression and cross-modality summarization, as interaction histories grow over time.
Integrating multimodal understanding of handwriting, diagrams, and speech would further close the gap between AI-based and human tutoring.
Finally, formalizing the conditions under which proactive agency improves learning outcomes, rather than merely system capability, remains an important open question for the community.

\section*{Limitations}

Our interactive evaluation relies on LLM-powered student simulators and rubric-based assessors, which inherit an irreducible gap between controlled simulation and real learner behavior.
Large-scale validation with human students is an essential next step.
For reasons of experimental tractability and evaluation cost, this report concentrates quantitative study on the foundational tutoring pipeline rather than attempting to fully measure every higher-level modality and proactive behavior in a single benchmark.
The general reasoning evaluation in Table~\ref{tab:generalization} compares the full pipeline, including its tool suite, against bare backbone models; isolating the contribution of the pipeline structure alone from the contribution of tool augmentation requires further experimentation.
The multi-stage pipeline trades additional inference cost for stronger controllability and personalization, and the economics of this trade-off depend on the deployment context.
Evaluation of the cognitive extensions and TutorBot requires longitudinal and deployment-based studies that remain future work.
The TutorBench construction pipeline could be extended to finer-grained courses, longer trajectories, and more diverse populations.
While we articulate a design for context-unified proactive tutoring, rigorously measuring its long-term impact on learning outcomes remains an open frontier.
